\paragraph{Thí nghiệm và kết quả}
\subparagraph{Chọn $\lambda$ tốt nhất bằng $k$-fold}
Chọn tập $\lambda$ ứng viên là:
\begin{equation}
  \lambda \in \left\{10^{-10}, 10^{-12}, \ldots, 10^{-1}, 1, 10\right\}.
\end{equation}

\subparagraph{Kết quả}
Kết quả thực nghiệm trên cả Ridge Regression và Lasso Regression cho thấy giá
trị $\lambda$ tối ưu nằm ở biên nhỏ nhất của tập ứng viên, cụ thể là $\lambda
= 10^{-10}$. Điều này cho thấy mô hình gần như không cần đến regularization và
nghiệm thu được xấp xỉ nghiệm của bài toán hồi quy tuyến tính thông thường
(Ordinary Least Squares).

\subparagraph{Giải thích}
Hiện tượng lựa chọn $\lambda$ rất nhỏ có thể xuất phát từ nguyên nhân mô hình
không cần regularization. Dữ liệu có số lượng đặc trưng nhỏ, mức nhiễu thấp
hoặc không có hiện tượng đa cộng tuyến mạnh, dẫn đến việc mô hình tuyến tính
đã tổng quát hóa tốt. Khi đó, việc thêm regularization chỉ làm tăng bias mà
không cải thiện variance, nên giá trị $\lambda$ tối ưu tiến về 0.

\subparagraph{Regularization path}
\figurename~\ref{fig:ridge-regularization-path} và
\figurename~\ref{fig:lasso-regularization-path} thể hiện sự thay đổi của các
hệ số hồi quy của Ridge và Lasso khi hệ số điều chỉnh (regularization
coefficient) thay đổi trên thang logarit.Khi $\lambda$ lớn (bên phải đồ thị),
các hệ số bị co lại đáng kể về gần $0$ do tác động mạnh của regularization.
Ngược lại, khi $\lambda$ giảm dần (bên trái đồ thị), các hệ số tiến gần tới
nghiệm của bài toán hồi quy tuyến tính không regularization (Ordinary Least
Squares). Điều này thể hiện rõ qua việc các đường biểu diễn trở nên ổn định và
đạt giá trị lớn hơn khi $\lambda$ dần đến $0$. Do nghiệm của Lasso Regression
thưa nên khi $\lambda$ tăng dần (bên phải
\figurename~\ref{fig:lasso-regularization-path}) nhiều hệ số hồi quy bị ép
bằng 0.
\begin{figure}[htbp]
  \centering
  \includegraphics[width=0.9\textwidth]{ridge-regularization-path}
  \caption{Regularization path của Ridge Regression}
  \label{fig:ridge-regularization-path}
\end{figure}

\begin{figure}[htbp]
  \centering
  \includegraphics[width=0.9\textwidth]{lasso-regularization-path}
  \caption{Regularization path của Lasso Regression}
  \label{fig:lasso-regularization-path}
\end{figure}

\subparagraph{Vùng tối ưu Elastic Net}
\figurename~\ref{fig:elastic-net-optimal-region} thể hiện sự thay đổi của MSE
trung bình theo sự thay đổi của $\lambda_1$ và $\lambda_2$ trên thang logarit.
Có thể quan sát rằng vùng tối ưu nằm ở vùng mà $\lambda_1$ và $\lambda_2$ rất
nhỏ (góc dưới trái của đồ thị). Điều này cho thấy mô hình đạt hiệu năng tốt
khi hai hệ số điều chỉnh (regularization coefficient) tiến dần về $0$. Điều
này chỉ ra rằng trong bài toán hiện tại, việc áp dụng regularization không
mang lại cải thiện rõ rệt về khả năng tổng quát hóa của mô hình.

\begin{figure}[htbp]
  \centering
  \includegraphics[width=0.9\textwidth]{elastic-net-optimal-region}
  \caption{Vùng tối ưu của Elastic Net}
  \label{fig:elastic-net-optimal-region}
\end{figure}

\subparagraph{Lựa chọn đặc trưng}
Kết quả thực nghiệm cho thấy khi chọn đặc trưng cả Forward Stepwise Selection
và Backward Elimination bằng mô hình OLS đều cho ra cùng một tập đặc trưng
\textit{(MedInc, HouseAge, Latitude, Longitude)} và chọn đặc trưng dựa trên
lasso cho kết quả là \textit{(MedInc, HouseAge, AveRooms, AveBedrms, AveOccup,
Latitude, Longitude)}

\subparagraph{Giải thích}

Sự khác biệt giữa kết quả lựa chọn đặc trưng của OLS và Lasso xuất phát từ
cơ chế regularization của hai phương pháp.

Đối với OLS, mô hình không áp dụng regularization nên việc lựa chọn đặc
trưng hoàn toàn dựa trên khả năng giải thích phương sai của từng biến.
Trong trường hợp này, cả Forward Stepwise Selection và Backward Elimination
đều hội tụ về cùng một tập đặc trưng \textit{(MedInc, HouseAge, Latitude,
Longitude)}, cho thấy đây là các biến mang thông tin quan trọng và ít dư
thừa nhất trong dữ liệu.

Ngược lại, Lasso sử dụng chuẩn $\ell_1$ để regularize các hệ số, từ đó có
khả năng đưa một số hệ số về 0 và thực hiện lựa chọn đặc trưng một cách tự
động. Tuy nhiên, trong thực nghiệm này, tập đặc trưng được chọn bởi Lasso
vẫn bao gồm nhiều biến, cụ thể là \textit{(MedInc, HouseAge, AveRooms,
AveBedrms, AveOccup, Latitude, Longitude)}. Điều này cho thấy giá trị
$\lambda$ tối ưu được lựa chọn tương đối nhỏ, khiến tác động của
regularization chưa đủ mạnh để loại bỏ hoàn toàn các đặc trưng có đóng góp
nhỏ.
